% !TeX root = ../main.tex

\acknowledgements

\iffalse
漫长的博士生涯接近尾声。在研究生学习期间，我得到了许多老师、同学以及朋友们的关心和帮助。在此，我想表达对他们的最深的感谢。
我虽未踏出过地球半步，却时常在理论物理浩瀚的宇宙中翱翔，俯瞰学术的广阔天空。

首先，衷心感谢我的指导教师，张本威教授。您是我报考华中师范大学时第一个接触的老师。
硕博期间，我在您的的悉心指导和教诲中成长。您言传身教，以深厚学识以及严谨的科研态度引领我步入高能核物理的殿堂，培养了我探索未知边界的本领。
您深刻的洞察力和智慧在许多方面都给予了我启发，让我学习如何独立做科学研究。
感谢您为我们提供了丰富的学术资源和良好的学术氛围，使我能够在国内外专业会议上登高而立，启迪思维，促进了我博士学业的发展。
您也常教导我们珍惜寸寸光阴，养成良好习惯。无论学术发展还是个人成长，您都对我产生了积极深远的影响。

% 张清
%感谢王恩科老师，每次在组会上总能一针见血的指出我们工作中的疏漏和不足，抓住被我们忽略的细节，是我们的研究更加深刻，对事物的认识更加全面。


感谢王新年老师，王恩科老师，蔡勖老师，张汉中老师，秦广友老师，侯德富老师，刘峰老师，肖博文老师，陈绍龙老师，付菁华老师，池丽平老师，郭红老师，施梳苏老师，毛亚显老师，伊航老师，计晨老师，王秋格老师的精彩授课和悉心指导。

感谢茹芃师兄在科研道路上对我的的耐心引导和大力帮助。您提供了许多启发性的宝贵建议并进行数次讨论，帮助我克服重重难关，照亮我前进的方向，重新点燃我对科研的信心，使得我的研究能够顺利推进。您的深耕细作以及高尚品格，让人钦佩，使我勇于迎接挑战并取得成果。
感谢代巍师兄，陈时勇师兄，你们慷慨的分享科研所得，让我在讨论和学习中获益良多。
感谢李汉林，何云存，申可明，江泽方，马国扬，嫣君，王洒，张善良，解曼，谢鑫，杜倩倩，康锦文，王磊，张清，张贺霞，罗愿，周琦，李瑶，孔唯玺，杨梦权，陶石磊，安敏，夏敏，陶俊琦，李小刚，樊翔，刘玉洁等师兄师姐师弟师妹们的无私帮助。我们一起探索科学的奥秘，共同奋斗，共同进步。
感谢粒子所的其它小伙伴李胜泰，高湛，张瑜，马飞，朱力强，Yasuki，Akio。我们一路陪伴，分享喜悦，并肩作战。感谢你们使我在科研道路之余感受到了无尽温暖，愿友谊长存！
感谢室友程瑄，王真真和古力姐，我们快乐相伴，共同度过了无数个美好日夜。

感谢我的外公外婆和奶奶，星空璀璨，是你们给予的温暖和关爱。你们是我前进的力量源泉。
感谢父母的无私关爱和辛勤付出，你们是最伟大的父母。感谢表弟表妹们、舅舅、姨等每一位亲人的支持。你们是我最坚实的后盾，是我的力量源泉，让我敢于追逐学术梦想并脚踏实地。
感谢我美丽故乡杜康酒厂，快乐记忆烙印心中，在成长的摇篮里塑造了我乐观向上的底色。

感谢馨香扑鼻的桂子山，卓越的教育资源和学术环境是我智慧的源泉。
感谢有着悠久历史和丰厚文化底蕴的祖国，感谢这个美好时代。在学习与成长的道路上，您为我提供了每一次宝贵机会，塑造了我的品格和思维方式。

博士之旅将告一段落，再次感谢你们的关怀和帮助，使我在面对困难时不屈不挠，我才能完成这段学术旅程。
希望你们能永远健康快乐。

% 以简短的文字表达作者对完成论文和学业提供帮助的\blind{老师}、\blind{同学}、\blind{亲属}等的感激之情。
\fi

\begin{signature}
  \blind{申淑婉} \\
%  \today \\
  2023年6月5日于桂子山
\end{signature}